\paragraph{Evolutionary dynamics as a nonlinear system.}
The formal setting is the standard one of evolutionary game
dynamics~\citep{taylor1978ess,hofbauer1998evolutionary,hofbauer2003evolutionary},
in which a population of competing designs is a nonlinear flow on a simplex,
qualitative changes of behaviour are bifurcations of that flow, and the objects
of interest are collective rather than
individual~\citep{perc2010coevolutionary,perc2017statphys}. That programme has
established across a wide range of social dilemmas that the transitions between
cooperative and defecting phases are sharp, that their location depends on
interaction structure at least as much as on payoffs, and that a dilemma can
often be reparameterised into a universal
form~\citep{perc2016phase,wang2015universal}; it has also mapped in detail how
much of the outcome is decided by who meets whom, through networks, layers or
conformity~\citep{nowak1992spatial,santos2005scale,ohtsuki2006simple,szabo2007evolutionary,wang2015multilayer,szolnoki2015conformity}.
Tags in particular have been studied as a way of generating assortment without
memory, both as phenotypic similarity in well-mixed
populations~\citep{antal2009phenotypic} and as a mutating tag space in finite
ones~\citep{traulsen2007chromodynamics}.

Two features separate the system studied here from that body of work. The first
is that the tag is forgeable at a rate the modeller sets, so the bifurcation
parameter is a property of a verification technology rather than of the
players, and the interesting comparative statics are with respect to an
engineering budget. The second is that the flow carries two attractors which
differ in whether the tag exists at all, rather than in how much cooperation it
supports, so the quantity that matters is the basin structure and not the
equilibrium level. Both features change what a policy conclusion can look like,
and neither is available in a model whose tag is inherited faithfully.

